\begin{thema}{Β}
Δίνεται η συνάρτηση $f(x) = \dfrac{x^2+a}{x-\beta}$.
\begin{erwthma}
\item Γνωρίζουμε ότι η ευθεία $x=1$ αποτελεί κατακόρυφη ασύμπτωτη της γραφικής παράστασης της $f$. Αυτό σημαίνει ότι ο παρονομαστής πρέπει να μηδενίζεται στο $x=1$ (και ο αριθμητής όχι). Eπομένως, $1 - \beta = 0 \Leftrightarrow \beta = 1$.
Η συνάρτηση γίνεται $f(x) = \dfrac{x^2+a}{x-1}$.
Στη συνέχεια, εκμεταλλευόμαστε ότι διέρχεται από το σημείο $A(2, 5)$, άρα $f(2) = 5$:
\[ \frac{2^2+a}{2-1} = 5 \Leftrightarrow \frac{4+a}{1} = 5 \Leftrightarrow a = 5 - 4 \Leftrightarrow a = 1 \]

\item Ο τύπος της συνάρτησης είναι $f(x) = \dfrac{x^2+1}{x-1}$ με πεδίο ορισμού το $D_f = \mathbb{R} \setminus \{1\}$.
Εφόσον ο βαθμός του αριθμητή (2) είναι ακριβώς κατά 1 μεγαλύτερος από τον βαθμό του παρανομαστή (1), η συνάρτηση διαθέτει \textbf{πλάγια ασύμπτωτη} (άρα όχι οριζόντια).
Κάνουμε την πολυωνυμική διαίρεση $(x^2+1) : (x-1)$ ή το σπάμε:
\[ f(x) = \frac{x^2-1+2}{x-1} = \frac{(x-1)(x+1)}{x-1} + \frac{2}{x-1} = x + 1 + \frac{2}{x-1} \]
Επειδή $\lim_{x\to\pm\infty} \frac{2}{x-1} = 0$, συνάγεται άμεσα ότι η ευθεία $y = x + 1$ είναι πλάγια ασύμπτωτη της $C_f$ τόσο στο $+\infty$ όσο και στο $-\infty$.

\item Η $f$ είναι παραγωγίσιμη στο πεδίο ορισμού της με:
\[ f'(x) = \frac{(x^2+1)'(x-1) - (x^2+1)(x-1)'}{(x-1)^2} = \frac{2x(x-1) - (x^2+1)}{(x-1)^2} = \frac{2x^2 - 2x - x^2 - 1}{(x-1)^2} = \frac{x^2-2x-1}{(x-1)^2} \]
Οι ρίζες της $f'(x)$ είναι οι λύσεις της δευτεροβάθμιας $x^2-2x-1=0$.
$\Delta = (-2)^2 - 4(1)(-1) = 4 + 4 = 8$.
Οι ρίζες είναι $x_{1,2} = \frac{2 \pm \sqrt{8}}{2} = \frac{2 \pm 2\sqrt{2}}{2} = 1 \pm \sqrt{2}$.
Το πρόσημο εξαρτάται μόνο από τον αριθμητή, αφού ο παρονομαστής είναι $(x-1)^2 > 0$.
\begin{center}
\begin{tikzpicture}
\tkzTabInit[color,lgt=2,espcl=1.5,colorC = red7,colorL = red9,colorV = red7]%
{$x$ / .8,$f'(x)$ /0.8,$f(x)$/0.8}%
{$-\infty$,$1-\sqrt{2}$,$1$,$1+\sqrt{2}$,$+\infty$}%
\tkzTabLine{ , +, z, -, d, -, z, +, }
\tkzTabLine{ , \nearrow, t, \searrow, d, \searrow, t, \nearrow, }
\end{tikzpicture}
\end{center}
Η συνάρτηση παρουσιάζει:
- \textbf{Τοπικό μέγιστο} στο $x = 1-\sqrt{2}$, με τιμή $f(1-\sqrt{2}) = \frac{(1-\sqrt{2})^2+1}{1-\sqrt{2}-1} = \frac{1-2\sqrt{2}+2+1}{-\sqrt{2}} = \frac{4-2\sqrt{2}}{-\sqrt{2}} = 2 - 2\sqrt{2}$.
- \textbf{Τοπικό ελάχιστο} στο $x = 1+\sqrt{2}$, με τιμή $f(1+\sqrt{2}) = \frac{(1+\sqrt{2})^2+1}{1+\sqrt{2}-1} = \frac{4+2\sqrt{2}}{\sqrt{2}} = 2 + 2\sqrt{2}$.

\item Η πλάγια ασύμπτωτη είναι η ευθεία $(\epsilon): y = x+1$.
Το ζητούμενο εμβαδόν μεταξύ της $C_f$, της ασύμπτωτης και των ευθειών $x=2$ και $x=e+1$ ισούται με το ολοκλήρωμα:
\[ E = \int_2^{e+1} |f(x) - (x+1)| dx = \int_2^{e+1} \left| \left( x+1+\frac{2}{x-1} \right) - (x+1) \right| dx = \int_2^{e+1} \left| \frac{2}{x-1} \right| dx \]
Για $x \in [2, e+1]$, έχουμε $x-1 > 0$, οπότε το απόλυτο παραλείπεται:
\[ E = \int_2^{e+1} \frac{2}{x-1} dx = \Big[ 2\ln|x-1| \Big]_2^{e+1} = 2\ln(e+1-1) - 2\ln(2-1) = 2\ln e - 2\ln 1 = 2 \cdot 1 - 0 = 2 \text{ τ.μ.} \]
\end{erwthma}
\end{thema}
